\section{Requerimientos funcionales}
% caracter ́ısticas, acciones, actores o entidades que definen de manera unica el flujo de trabajo de la organizaci ́on.

\subsection*{Características}
El prototipo desarrollado debe proporcionar una forma de solicitar, almacenar y manejar la información necesaria en archivos .csv

\subsection*{Entidades y actores}
Las entidades que nos permitirán el manejo de dicha información serán:
\begin{itemize}
    \item Sucursal: con un identificador entero y único, nombre, dirección y teléfono.
    \item Premio: con un identificador entero y único, nombre, descripción, puntos requeridos y existencias.
    \item Cliente: con un identificador entero y único, nombre, correo, teléfono y puntos acumulados.
\end{itemize}

Los actores principales que estarán en contacto con el prototipo serán:
\begin{itemize}
    \item Administrador del centro: Quien será el usuario final del prototipo y quien administra la información del centro.
\end{itemize}

\subsection*{Acciones/operaciones CRUD}
El manejo de la información sobre sucursales, premios y clientes permitirá las siguientes acciones:
\begin{itemize}
    \item Agregar información sobre sucursales, premios y clientes.
    \item Consultar la información sobre las sucursales, premios y clientes.
    \item Editar información sobre sucursales, premios y clientes.
    \item Eliminar información sobre sucursales, premios y clientes.
    \item Guardar la información sobre estas tres entidades en archivos .csv de forma persistente.
    \item Leer la información de estos archivos para poder realizar las consultas en cuqluier momento.
    \item Organizar la información sobre los tres entes de forma adecuada.
\end{itemize}

\subsection*{Interfaz de usuario}
Se proporcionará un menú de interacción que permite realizar las operaciones antes descritas para cada entidad.

\subsection*{Consultas con llave}
El sistema puede recibir una llave de entidad para devolver toda la información relacionada.

\subsection*{Validación}
El prototipo puede almacenar únicamente el formato específicado para cada uno de los parámetros definidos de cada entidad.

\subsection*{Manejo de excepciones}
El usuario debe poder observar el error cometido durante la ejecución del prototipo.






